% !TEX root = ../main.tex
\section{Dados e Métodos}
\label{sec:dados}

\subsection{Proveniência dos dados}

Duas fontes brutas independentes alimentam o pipeline:

\begin{itemize}
  \item \textbf{LiDAR aerotransportado} --- nuvens de pontos \texttt{.laz} do
    programa \textit{NASA Carbon Monitoring System (CMS) Brazil}, distribuídas
    pelo ORNL DAAC. O inventário completo de tiles disponibilizado pela missão
    soma \textbf{3.152 tiles}, cobrindo múltiplas campanhas de voo (diferentes
    anos e regiões) sobre a Amazônia brasileira.
  \item \textbf{Inventário florestal de campo} --- arquivos \texttt{.csv} com
    medições árvore a árvore (diâmetro à altura do peito, altura, espécie) e
    \texttt{.kmz} com a geometria das parcelas, cobrindo \textbf{28 sites} de
    inventário distribuídos pela Amazônia.
\end{itemize}

As duas fontes são coletadas, processadas e cruzadas de forma independente uma da
outra até a etapa de interseção espacial --- nenhuma delas é derivada da outra.

\subsection{Pipeline de pré-processamento}

O cruzamento entre as duas fontes segue seis etapas, ilustradas na
Figura~\ref{fig:pipeline}.

\begin{figure}[H]
  \centering
  % PLACEHOLDER — substituir pela figura do pipeline:
  % \includegraphics[width=0.85\textwidth]{figura1_pipeline.png}
  \figplaceholder{[PLACEHOLDER] Figura 1 --- Pipeline de pré-processamento (seis
  etapas): extração de geometrias, interseção espacial parcela $\times$ tile,
  recorte do LiDAR por parcela, normalização do inventário, cálculo de biomassa e
  pareamento final.}
  \caption{Pipeline de pré-processamento que cruza o LiDAR aerotransportado com o
  inventário florestal de campo.}
  \label{fig:pipeline}
  \fonte{elaborado pelo autor (2026)}
\end{figure}

A definição de \textbf{parcela} (\texttt{plot\_id}) merece nota metodológica:
auditorias sobre os KMLs de inventário revelaram parcelas com o campo
\texttt{Name} não único --- polígonos distintos (quadras) compartilhando o mesmo
nome ---, o que levava o pipeline original a dissolver quadras diferentes numa só,
colapsando contagens e pareando AGB de campo com o clip de LiDAR errado. A
correção (documentada \textit{in extenso} no apêndice de qualidade de dados)
define o \texttt{plot\_id} por uma regra hierárquica: nome único $\rightarrow$ usa
o nome; nome repetido $\rightarrow$ sufixo sequencial (\texttt{T20\_1},
\texttt{T20\_2}, \dots); nome não informativo $\rightarrow$ índice sequencial.
Essa correção, mais a remoção de feições geometricamente duplicadas e a
recuperação de parcelas indevidamente descartadas por bugs de leitura de
coordenadas, elevou o conjunto de parcelas treináveis de 383 para as 493 finais
--- um ganho de 29\% no tamanho da amostra de treino, obtido inteiramente por
correção de bugs, sem novos dados brutos.

O critério de qualificação espacial usa \textbf{$\geq$ 99,9\% da área da parcela
dentro da cobertura contígua de uma campanha} (tiles vizinhos do mesmo voo,
unidos; subcampanhas do mesmo ano tratadas como uma só), em vez de um teste
\texttt{within} estrito --- o que evitava que parcelas 100\% cobertas, mas tocando
a borda da união de tiles, fossem indevidamente descartadas.

\subsection{Números do dataset}

\begin{table}[H]
  \centering
  \caption{Números do conjunto de dados ao longo do pipeline.}
  \label{tab:dataset}
  \begin{tabular}{lr}
    \toprule
    \textbf{Métrica} & \textbf{Valor} \\
    \midrule
    Tiles no inventário LiDAR da NASA & 3.152 \\
    Sites de inventário (total) & 28 \\
    Tiles LiDAR distintos efetivamente usados & 242 \\
    Parcelas no inventário (uma por quadra, pós-deduplicação) & 581 \\
    Pares parcela $\times$ tile & 1.287 \\
    Parcelas com cobertura LiDAR ($\geq$ 99,9\% da área) & 553 \\
    Parcelas com AGB de campo calculável (nível quadra) & 509 \\
    \textbf{Parcelas treináveis --- AGB $\times$ LiDAR pareados} & \textbf{493} \\
    Parcelas treináveis após remoção de outliers geométricos & 325 \\
    \bottomrule
  \end{tabular}
  \fonte{auditoria de dados do projeto (última varredura
  registrada), reproduzida integralmente na página de Apêndice --- Qualidade de
  Dados}
\end{table}

\subsection{Biomassa de referência (alvo)}
\label{sec:biomassa}

A AGB de campo é calculada árvore a árvore e agregada por parcela, sob três
fórmulas alométricas candidatas que diferem na especificidade da densidade de
madeira ($\rho$) e no tratamento de árvores mortas e palmeiras:

\begin{table}[H]
  \centering
  \caption{Fórmulas alométricas candidatas para a biomassa de referência.}
  \label{tab:formulas}
  \begin{tabular}{p{3.2cm}p{6.2cm}p{4.6cm}}
    \toprule
    \textbf{Fórmula} & \textbf{Equações usadas} & \textbf{Densidade de madeira ($\rho$)} \\
    \midrule
    \textbf{M1 --- Chave uniforme} & Chave et al. (2014) para todas as árvores &
      $\rho$ = 0,6 g/cm\textsuperscript{3} (média global) \\
    \textbf{M2 --- Chave por tipo} & Chave (2014, vivas) $\cdot$ Chambers (2000,
      mortas) $\cdot$ Goodman (2013, palmeiras) & $\rho$ = 0,6 g/cm\textsuperscript{3}
      (média global) \\
    \textbf{M3 --- Chave por tipo + $\rho$ específica} & Mesmas equações de M2 &
      $\rho$ por espécie (GWDD Zanne/Chave; \textit{fallback} gênero
      $\rightarrow$ site $\rightarrow$ 0,6) \\
    \bottomrule
  \end{tabular}
  \fonte{elaborado pelo autor (2026), com base em Chave et al. (2014), Chambers
  et al. (2000) e Goodman et al. (2013)}
\end{table}

Equações de carbono individual (fração de carbono $f_C = 0,5$):

\begin{align}
  \text{IAGC}_{\text{viva}}     &= 0{,}0673 \times 0{,}5 \times
     \left(\rho \cdot \text{DBH}^2 \cdot H\right)^{0{,}976}
     && \text{--- Chave et al. (2014), Eq.~7} \\
  \text{IAGC}_{\text{morta}}    &= 0{,}1007 \times 0{,}5 \times 0{,}40 \times
     \text{DBH}^2 \cdot H^{0{,}818}
     && \text{--- Chambers et al. (2000)} \\
  \text{IAGC}_{\text{palmeira}} &= 0{,}03781 \times 0{,}5 \times
     \text{DBH}^{2{,}7483}
     && \text{--- Goodman et al. (2013)}
\end{align}

com $\text{AGB} = 2 \times \text{AGC}$ (unidade final: Mg/ha). \textbf{Este
trabalho usa M1 como alvo de treino para todos os modelos.} Essa escolha --- e o
fato de não haver, nesta versão, uma análise de sensibilidade dos resultados a
M2/M3 --- é revisitada criticamente na Seção~\ref{sec:allom-sens}.

Adicionalmente, corrige-se o \textbf{descompasso temporal} entre o ano do voo
LiDAR e o ano da medição de campo: o DBH é ``crescido'' até o ano do voo pelo
incremento anual observado (por remedição, quando disponível; senão pela mediana
do site; senão um valor \textit{default}), e altura e biomassa são re-derivadas a
partir do DBH ajustado. O efeito médio é pequeno ($\sim$1,3\% na biomassa média),
mas chega a 7--12\% nos sites com gap de 2--3 anos entre voo e medição --- não
desprezível para parcelas individuais, mesmo que pequeno em média.

\subsection{Abordagens de modelagem comparadas}

Seis representações do sinal LiDAR foram comparadas, da mais compacta à mais rica
em graus de liberdade --- todas treinadas sobre o \textbf{mesmo conjunto de
parcelas}, com o mesmo protocolo de validação cruzada, para que a diferença de
desempenho seja atribuível à representação e ao modelo, não ao dado de treino:

\begin{table}[H]
  \centering
  \caption{Abordagens de modelagem comparadas.}
  \label{tab:modelos}
  \small
  \begin{tabular}{clp{5.2cm}p{3.2cm}}
    \toprule
    \textbf{\#} & \textbf{Abordagem} & \textbf{Entrada} & \textbf{Modelo} \\
    \midrule
    1 & GBR --- alturas ordenadas & 1024 alturas de dossel amostradas e ordenadas
      + log(n pontos) --- 1025 atributos & Gradient Boosting \\
    2 & GBR --- $K$ quantis & $K$ quantis igualmente espaçados da altura +
      log(n pontos) --- $K{+}1$ atributos, $K \in \{8,16,32,64,128,256\}$ &
      Gradient Boosting \\
    3 & ABA (\textit{Area-Based Approach}) & 37 métricas estruturais --- percentis
      de altura, densidade vertical, rugosidade convolucional do CHM &
      Gradient Boosting \\
    4 & Random Forest estrutural & As mesmas 37 métricas do ABA &
      Random Forest (\textit{bagging}) \\
    5 & PointNet & Nuvem de pontos crua $(x, y, \text{altura})$, até 8192
      pontos/parcela, 2048 amostrados por passo & PointNet (MLP compartilhado +
      max-pool) \\
    6 & CNN 2D --- raster & Grade 32$\times$32, 4 canais (altura máx., altura
      média, densidade log, desvio-padrão) & CNN 2D (3 blocos conv+BN+ReLU) \\
    7 & CNN 3D --- voxel & Volume 16$\times$16$\times$16, 2 canais (ocupação,
      densidade log) & CNN 3D (3 blocos conv3d+BN+ReLU) \\
    \bottomrule
  \end{tabular}
  \fonte{elaborado pelo autor (2026)}
\end{table}

As abordagens 3 e 4 compartilham exatamente as mesmas 37 features de entrada ---
a única diferença é \textit{bagging} (Random Forest) vs. \textit{boosting}
(Gradient Boosting) ---, isolando o efeito do algoritmo de ensemble do efeito da
representação. As abordagens 5--7 recebem representações geometricamente
crescentes em riqueza (nuvem crua $\rightarrow$ grade 2D $\rightarrow$ volume 3D)
e, com isso, em número de parâmetros e propensão a sobreajuste com poucas
amostras.

\subsection{Protocolo de validação}
\label{sec:protocolo}

Todos os modelos são avaliados por \textbf{validação cruzada k-fold aleatória
(5 folds)}, com o alvo transformado em log1p. ABA e Random Forest usam os mesmos
hiperparâmetros de Gradient Boosting (\texttt{n\_estimators=300}, \texttt{max\_depth=4},
\texttt{learning\_rate=0,05}, \texttt{subsample=0,8}) para que a comparação isole
o efeito do algoritmo. Reporta-se RMSE (Mg/ha), $R^2$ e rRMSE (RMSE relativo à
média observada, \%) --- agregados sobre as previsões \textit{out-of-fold} de
todos os folds (métrica ``global''), não a média simples das métricas por fold.

Cada modelo é avaliado em \textbf{seis variantes do conjunto de treino, cada uma
testando UM tratamento isoladamente contra a mesma nuvem crua} --- não empilhadas
umas sobre as outras, para que o efeito de cada tratamento possa ser atribuído a
ele mesmo, não à combinação (Seção~\ref{sec:preproc}): \textbf{sem processamento}
(\texttt{raw}, baseline), \textbf{só projeção do solo} (\texttt{ground\_only} ---
relevo local célula-a-célula, sem filtro), \textbf{só filtro $\leq$1,3~m}
(\texttt{filter\_only} --- remoção dos retornos até a altura do peito, aplicada
sobre o mesmo deslocamento mínimo de Z do baseline cru, sem acionar a projeção de
relevo completa), \textbf{só remoção de outliers} (\texttt{outliers\_only} ---
parcelas com área $>$ 0,5~ha ou compacidade $<$ 0,6 removidas, sem nenhum
tratamento na nuvem), \textbf{só correção de gap temporal} (\texttt{gap\_only} ---
alvo realinhado à data do voo LiDAR, sem nenhum tratamento na nuvem) e
\textbf{tudo combinado} (\texttt{gap\_corrected\_noout} --- os quatro tratamentos
juntos, o pipeline completo). Variantes intermediárias empilhadas (solo+filtro,
+outliers, +gap) também foram treinadas e permanecem disponíveis nos dados, mas a
análise principal da Seção~\ref{sec:preproc} usa as seis variantes isoladas como
evidência primária, por permitirem atribuir o ganho de desempenho a um tratamento
específico em vez de a uma combinação.

\begin{quote}
\textbf{Nota de escopo (ver Seção~\ref{sec:spatial-validation}):} o k-fold aqui é
aleatório sobre parcelas, não sobre sítios --- parcelas do mesmo site de
inventário podem cair em treino e teste no mesmo fold. Isso é adequado para medir
o ajuste do modelo à distribuição de estrutura do dossel presente na amostra, mas
\textbf{infla a métrica de generalização espacial} por autocorrelação entre
parcelas vizinhas do mesmo sítio --- a leitura correta dos números desta seção é
``quanto da variância intra-amostra é explicada'', não ``quão bem isto generaliza
para floresta amazônica não amostrada''.
\end{quote}
